\documentclass[11pt,a4paper]{article}
\usepackage[T1]{fontenc}
\usepackage[utf8]{inputenc}
\usepackage{lmodern}
\usepackage[margin=1.1in]{geometry}
\usepackage{booktabs}
\usepackage{microtype}
\usepackage{enumitem}
\usepackage[hidelinks]{hyperref}

\title{Stopping the stop: a self-labeling gate for premature termination\\
       in autonomous coding agents}
\author{Juan Lescano}
\date{29 August 2026}

\begin{document}
\maketitle

\begin{abstract}
A coding agent granted standing permission still ends its turn early: it names
the next step instead of taking it, offers a menu instead of choosing, or hands
the user a command it could have run. We describe a gate on the agent's stop
path that detects this in three lanes of increasing cost -- an audit of the
agent's own escape hatch, 298 pattern labeling functions, and a 9B local judge
-- and re-wakes the agent when the closing message shows the pattern. Building
the gate was straightforward; tuning it was not, because tuning needs labels and
labels need a reader. We show that two usable labels are already present in the
transcript and cost no annotation. Before the gate exists, the user's own next
message supplies a weak label for premature stopping. After the gate
intervenes, the agent's tool calls before the user speaks again supply a label
for the intervention itself. We evaluate over 854 closing pairs from 66 real
sessions, ablate each lane against the first label, and describe a loop that
uses the second to name its own patterns for demotion. We report negative
results in full, including two failed attempts to obtain a usable confidence
signal from the judge, and one case where an exception a 9B model could not hold
in its prompt became trivial once moved out of the model entirely.
\end{abstract}
